% ============================================================
%  CUADERNILLO DE ESTUDIO — CLASE 9
%  Seguridad de la Información y de Sistemas
%  Lic. en Gestión de Negocios Digitales — UCA
%  Compilar con: pdflatex (dos pasadas para el índice)
% ============================================================

\documentclass[12pt,a4paper]{article}

% ---------- Paquetes ----------
\usepackage[utf8]{inputenc}
\usepackage[spanish]{babel}
\usepackage[T1]{fontenc}
\usepackage{lmodern}
\usepackage[margin=2.2cm,headheight=14pt]{geometry}
\usepackage{amsmath}
\usepackage{amssymb}
\usepackage{enumitem}
\usepackage{booktabs}
\usepackage{array}
\usepackage{tabularx}
\usepackage[table]{xcolor}
\usepackage{tcolorbox}
\usepackage{graphicx}
\usepackage{fancyhdr}
\usepackage{titlesec}
\usepackage{hyperref}
\usepackage{multicol}
\usepackage{pifont}
\usepackage{wasysym}
\usepackage{tikz}
\usetikzlibrary{positioning,arrows.meta,shapes.geometric}

% ---------- Colores ----------
\definecolor{azulUCA}{HTML}{003366}
\definecolor{doradoUCA}{HTML}{B8860B}
\definecolor{grisClaro}{HTML}{F2F2F2}
\definecolor{rojoAlerta}{HTML}{C0392B}
\definecolor{verdeOK}{HTML}{27AE60}
\definecolor{azulCielo}{HTML}{2980B9}
\definecolor{naranjaAmenaza}{HTML}{D35400}
\definecolor{violetaIA}{HTML}{8E44AD}
\definecolor{azulLegal}{HTML}{154360}
\definecolor{moradoEtica}{HTML}{6C3483}
\definecolor{azulCookie}{HTML}{2874A6}

% ---------- tcolorbox estilos ----------
\tcbuselibrary{skins,breakable}

\newtcolorbox{cajaTitulo}[1][]{
  colback=azulUCA, colframe=azulUCA, coltext=white,
  fonttitle=\Large\bfseries, arc=4mm, #1
}

\newtcolorbox{cajaDefinicion}{
  colback=azulCielo!10, colframe=azulCielo,
  fonttitle=\bfseries, title=Definici\'on, arc=3mm, breakable
}

\newtcolorbox{cajaLegal}{
  colback=azulLegal!8, colframe=azulLegal,
  fonttitle=\bfseries\color{azulLegal},
  title={\ding{211}\; Marco Legal}, arc=3mm, breakable
}

\newtcolorbox{cajaEtica}{
  colback=moradoEtica!8, colframe=moradoEtica,
  fonttitle=\bfseries\color{moradoEtica},
  title={\ding{73}\; \'Etica y Privacidad}, arc=3mm, breakable
}

\newtcolorbox{cajaCookie}{
  colback=azulCookie!8, colframe=azulCookie,
  fonttitle=\bfseries\color{azulCookie},
  title={\ding{80}\; Cookies y Rastreo}, arc=3mm, breakable
}

\newtcolorbox{cajaActividad}{
  colback=grisClaro, colframe=azulUCA,
  fonttitle=\bfseries\color{azulUCA},
  title={\ding{45}\; Actividad Pr\'actica}, arc=3mm, breakable
}

\newtcolorbox{cajaNota}{
  colback=rojoAlerta!8, colframe=rojoAlerta,
  fonttitle=\bfseries\color{rojoAlerta}, title=Atenci\'on Negocio, arc=3mm, breakable
}

% ---------- Header / Footer ----------
\pagestyle{fancy}
\fancyhf{}
\fancyhead[L]{\small\textcolor{azulUCA}{Seguridad de la Informaci\'on y de Sistemas}}
\fancyhead[R]{\small\textcolor{azulUCA}{Cuadernillo --- Clase 9}}
\fancyfoot[C]{\thepage}
\fancyfoot[L]{\small\textcolor{gray}{UCA}}
\renewcommand{\headrulewidth}{0.4pt}
\renewcommand{\footrulewidth}{0.4pt}

% ---------- Títulos de sección ----------
\titleformat{\section}
  {\color{azulUCA}\Large\bfseries}
  {\thesection.}{0.6em}{}
\titleformat{\subsection}
  {\color{azulCielo}\large\bfseries}
  {\thesubsection}{0.6em}{}

% ---------- Enlaces ----------
\hypersetup{colorlinks=true, linkcolor=azulUCA, urlcolor=azulCielo}

% ---------- Checkbox ----------
\newcommand{\checkboxVacio}{$\square$\;}
\newcommand{\checkMark}{\ding{51}}

% ============================================================
\begin{document}

% ============================================================
%  PORTADA
% ============================================================
\begin{titlepage}
\centering
\vspace*{1cm}

\begin{cajaTitulo}[width=\textwidth, halign=center]
  SEGURIDAD DE LA INFORMACI\'ON Y DE SISTEMAS\\[4pt]
  {\large Licenciatura en Gesti\'on de Negocios Digitales --- UCA}
\end{cajaTitulo}

\vspace{1.5cm}

{\Huge\bfseries\color{azulUCA} CUADERNILLO DE ESTUDIO}\\[8pt]
{\LARGE\color{doradoUCA} CLASE 9}\\[12pt]
{\Large\itshape ``Marco legal II: cibercrimen, \'etica,\\[2pt] privacidad por dise\~no y cookies''}\\[6pt]
{\large\color{azulCookie}\ding{80}\; Cierre de la Unidad 3}\\[1.5cm]

\begin{tabular}{rl}
\textbf{Profesor:}   & Mg.\ Ing.\ Rodrigo Atilio Elgueta \\[4pt]
\textbf{Unidad:}     & Cierre de Unidad 3 \\[4pt]
\end{tabular}

\vfill

\begin{tcolorbox}[colback=grisClaro, colframe=azulUCA, arc=3mm, width=0.9\textwidth]
  \centering
  {\bfseries Datos del/la estudiante}\\[10pt]
  Nombre y apellido: \underline{\hspace{5cm}}\\[10pt]
  Grupo N.\textdegree: \underline{\hspace{5cm}}\\[10pt]
  Negocio Digital (TP): \underline{\hspace{5cm}}\\[10pt]
  Fecha: \underline{\hspace{5cm}}
\end{tcolorbox}

\vspace{1cm}
{\small\textcolor{gray}{Material de uso exclusivo para la cursada.}}
\end{titlepage}

% ============================================================
%  ÍNDICE
% ============================================================
\tableofcontents
\newpage

% ============================================================
%  SECCIÓN 1 — INTRODUCCIÓN
% ============================================================
\section{De la protecci\'on al delito: las reglas del juego}

En la Clase 8 vimos c\'omo la ley (Ley 25.326 y GDPR) protege los datos de los usuarios y les otorga derechos (ARCO). Pero, ?qu\'e pasa cuando alguien cruza la l\'inea y ataca esos datos? ?Y qu\'e pasa cuando una empresa, sin llegar a cometer un delito, dise\~na su producto de forma manipuladora para extraer la mayor cantidad de informaci\'on posible?

Hoy cerramos la \textbf{Unidad 3} explorando tres dimensiones fundamentales para cualquier negocio digital:
\begin{enumerate}[leftmargin=2em]
    \item \textbf{El Cibercrimen:} Qu\'e conductas est\'an tipificadas como delitos inform\'aticos en Argentina.
    \item \textbf{La \'Etica y Privacidad por Dise\~no:} C\'omo construir productos que respeten al usuario desde su arquitectura, no solo como un parche legal.
    \item \textbf{Las Cookies:} El mecanismo t\'ecnico y legal que sostiene la publicidad digital y el rastreo en la web (uno de los temas m\'as pedidos por ustedes en la encuesta).
\end{enumerate}

\newpage
% ============================================================
%  SECCIÓN 2 — LEY 26.388 (DELITOS INFORMÁTICOS)
% ============================================================
\section{Cibercrimen en Argentina: Ley 26.388}

Sancionada en 2008, la Ley 26.388 modific\'o el C\'odigo Penal argentino para incorporar los ``delitos inform\'aticos''. Antes de esta ley, atacar un servidor o robar datos digitales era un vac\'io legal o se juzgaba con figuras anal\'ogicas (como ``robo'' o ``da\~no''), que no encajaban bien con la naturaleza intangible de los datos.

\subsection{Principales figuras penales incorporadas}
\renewcommand{\arraystretch}{1.5}
\begin{tabularx}{\textwidth}{|>{\bfseries\color{azulLegal}}p{4cm}|X|}
\hline
\rowcolor{grisClaro}
\textbf{Delito} & \textbf{?Qu\'e sanciona?} \\
\hline
Acceso indebido & Entrar a un sistema o dato restringido sin autorizaci\'on (hacking b\'asico). \\
\hline
Retenci\'on de datos & Apoderarse de datos inform\'aticos o retenerlos contra la voluntad de su due\~no. \\
\hline
Da\~no inform\'atico & Alterar, destruir o inutilizar datos, documentos o sistemas (ej. borrar una base de datos, lanzar un wiper). \\
\hline
Fraude inform\'atico & Manipular datos para obtener un beneficio econ\'omico injusto (ej. alterar el precio en un e-commerce, clonar tarjetas). \\
\hline
Interceptaci\'on de comunicaciones & Captar comunicaciones electr\'onicas no destinadas al p\'ublico (ej. sniffing de redes WiFi abiertas). \\
\hline
Pornograf\'ia infantil & Producci\'on, tenencia o distribuci\'on de material de abuso sexual infantil a trav\'es de medios digitales. \\
\hline
Grooming & Contacto por medios tecnol\'ogicos de un adulto con un menor con el prop\'osito de cometer un delito contra su integridad sexual. \\
\hline
\end{tabularx}

\begin{cajaNota}
\textbf{Para el negocio digital:} Conocer esta ley no es solo para no cometer delitos. Es vital para saber \textbf{cu\'ando un incidente de seguridad (como un hackeo o una extorsi\'on) debe ser denunciado a la fiscal\'ia} como un delito penal, en lugar de manejarse solo como un problema t\'ecnico interno.
\end{cajaNota}

\newpage
% ============================================================
%  SECCIÓN 3 — ÉTICA Y PRIVACIDAD POR DISEÑO
% ============================================================
\section{\'Etica y Privacidad por Dise\~no (Privacy by Design)}

Cumplir la ley es el m\'inimo exigible. Pero los negocios digitales que generan \textbf{confianza a largo plazo} van m\'as all\'a: aplican la \'etica y la Privacidad por Dise\~no.

\begin{cajaEtica}[title={?Qu\'e es Privacidad por Dise\~no?}]
Es un enfoque propuesto por Ann Cavoukian que establece que la privacidad no debe agregarse como un parche al final del desarrollo de un producto, sino que debe estar \textbf{incrustada en la arquitectura misma} del sistema desde el primer d\'ia.
\end{cajaEtica}

\subsection{Los 7 principios fundamentales}
\begin{enumerate}[leftmargin=2em]
    \item \textbf{Proactivo, no reactivo:} Anticipar y prevenir incidentes de privacidad antes de que ocurran.
    \item \textbf{Privacidad como configuraci\'on por defecto:} El usuario no deber\'ia tener que hacer clic en 20 men\'us para proteger sus datos. Si no hace nada, sus datos est\'an protegidos (Default Privacy).
    \item \textbf{Privacidad embebida en el dise\~no:} Es parte integral del producto, no un a\~nadido.
    \item \textbf{Funcionalidad total (Suma positiva):} No se trata de elegir entre seguridad y usabilidad; se buscan soluciones que ofrezcan ambas (Win-Win).
    \item \textbf{Seguridad de extremo a extremo:} Protecci\'on durante todo el ciclo de vida del dato, desde que se recolecta hasta que se destruye.
    \item \textbf{Visibilidad y transparencia:} Las pr\'acticas de datos deben ser verificables por los usuarios.
    \item \textbf{Respeto por la privacidad del usuario:} El inter\'es del usuario es la prioridad (empoderamiento a trav\'es de los derechos ARCO).
\end{enumerate}

\begin{cajaActividad}[title={Reflexi\'on \'Etica: Dark Patterns}]
Los \textbf{Dark Patterns} (patrones oscuros) son dise\~nos de interfaz enga\~nosos que manipulan al usuario para que haga algo que no quer\'ia hacer (ej. suscribirse a un servicio, aceptar cookies de marketing, o hacer que el bot\'on de ``Cancelar suscripci\'on'' sea casi invisible).\\[6pt]
\textbf{Pregunta para el grupo:} Si un Dark Pattern aumenta las ventas un 15\% a corto plazo pero es legal (o est\'a en un gris legal), ?deber\'ia usarlo un negocio digital? ?Qu\'e impacto tiene en la confianza de la marca a largo plazo?
\end{cajaActividad}

\newpage
% ============================================================
%  SECCIÓN 4 — COOKIES Y RASTREO
% ============================================================
\section{El universo de las Cookies}

El protocolo HTTP (la base de la web) es ``sin estado'' (stateless). Esto significa que el servidor no recuerda qui\'en sos de una p\'agina a otra. Las \textbf{cookies} se inventaron para resolver esto: son peque\~nos archivos de texto que el servidor le pide al navegador que guarde y le devuelva en cada visita.

\subsection{Tipos de Cookies}

\renewcommand{\arraystretch}{1.5}
\begin{tabularx}{\textwidth}{|>{\bfseries\color{azulCookie}}p{3.5cm}|X|X|}
\hline
\rowcolor{grisClaro}
\textbf{Tipo} & \textbf{Caracter\'istica} & \textbf{Ejemplo de uso} \\
\hline
\textbf{De Sesi\'on} & Se borran autom\'aticamente al cerrar el navegador. & Mantener el carrito de compras mientras naveg\'as por la tienda. \\
\hline
\textbf{Persistentes} & Tienen una fecha de expiraci\'on y quedan en tu disco. & Recordar tu idioma preferido o mantener tu sesi\'on iniciada (``Recordarme''). \\
\hline
\textbf{De Origen (First-party)} & Creadas por el dominio que est\'as visitando directamente. & Autenticaci\'on, preferencias del sitio. \\
\hline
\textbf{De Terceros (Third-party)} & Creadas por dominios distintos al que visitas (ej. anunciantes, trackers). & Publicidad comportamental, botones de ``Me gusta'' de redes sociales, anal\'itica cruzada. \\
\hline
\end{tabularx}

\begin{cajaCookie}[title={El problema del consentimiento y la ``fatiga de cookies''}]
Por normativas como el GDPR y la Ley 25.326, los sitios deben pedir consentimiento para cookies \textbf{no esenciales} (marketing, anal\'itica de terceros). Sin embargo, la mayor\'ia de los sitios usan Dark Patterns: el bot\'on de ``Aceptar Todo'' es gigante y verde, mientras que ``Rechazar'' o ``Configurar'' est\'a escondido en gris o requiere 5 clics. Esto viola el principio de \textbf{Privacidad por Defecto}.
\end{cajaCookie}

\newpage
% ============================================================
%  SECCIÓN 5 — ACTIVIDAD PRÁCTICA
% ============================================================
\section{Taller: Radiograf\'ia de Cookies}

Hoy vamos a dejar de ser usuarios pasivos y vamos a ``ver la matrix'' usando las Herramientas de Desarrollador (DevTools) del navegador.

\begin{cajaActividad}[title={Inspecci\'on en tiempo real}]
\textbf{Paso 1:} Abran una ventana de \textbf{Inc\'ognito} en Chrome, Edge o Firefox (para no arrastrar cookies viejas).\\
\textbf{Paso 2:} Ingresen a un portal de noticias muy concurrido o a un e-commerce grande (ej. clarin.com, mercadolibre.com).\\
\textbf{Paso 3:} \textbf{NO} hagan clic en el banner de cookies todav\'ia. Presionen \textbf{F12} (o clic derecho $\rightarrow$ Inspeccionar).\\
\textbf{Paso 4:} Vayan a la pesta\~na \textbf{Application} (o \textbf{Storage} en Firefox) $\rightarrow$ \textbf{Cookies}.\\
\textbf{Paso 5:} Respondan:
\end{cajaActividad}

\vspace{4pt}
\renewcommand{\arraystretch}{1.5}
\begin{tabularx}{\textwidth}{|X|}
\hline
\textbf{1. Sin haber aceptado nada, ?cu\'antas cookies ya se cargaron autom\'aticamente? ?De qu\'e dominios son?} \\[2.5cm]
\hline
\textbf{2. Busquen una cookie que parezca de rastreo (nombres como \_ga, \_fbp, fr, c\_user). ?Qu\'e dato est\'a almacenando?} \\[2.5cm]
\hline
\textbf{3. Ahora hagan clic en ``Rechazar todo'' o ``Configurar'' (si el sitio lo permite). ?Qu\'e cookies desaparecieron y cu\'ales se quedaron (las esenciales)?} \\[2.5cm]
\hline
\textbf{4. Reflexi\'on de Negocio: Si ustedes estuvieran dise\~nando el banner de cookies de su TP Integrador, ?c\'omo lo har\'ian para ser \'eticos y cumplir Privacidad por Defecto?} \\[3cm]
\hline
\end{tabularx}

\newpage
% ============================================================
%  SECCIÓN 6 — CUESTIONARIO 3
% ============================================================
\section{Cuestionario 3 de autoevaluaci\'on}

\begin{cajaNota}[title={Instrucciones}]
Este cuestionario cierra la Unidad 3. Se realiza al final de la clase (Kahoot / Forms / papel). \textbf{No es eliminatorio.} Sirve de autoevaluaci\'on y repaso para el Parcial 2 (Clase 11). Las respuestas correctas est\'an marcadas con \checkMark \; al final.
\end{cajaNota}

\vspace{6pt}

\textbf{1. La Ley 25.326 en Argentina regula:}
\begin{enumerate}[label=\alph*), leftmargin=2.5em]
  \item El comercio electr\'onico
  \item La protecci\'on de datos personales
  \item La propiedad intelectual
  \item Los delitos inform\'aticos exclusivamente
\end{enumerate}

\textbf{2. Los derechos ARCO permiten a una persona:}
\begin{enumerate}[label=\alph*), leftmargin=2.5em]
  \item Acceder, Rectificar, Cancelar (Suprimir) y Oponerse al tratamiento de sus datos
  \item Auditar, Reclamar, Comprar y Ofertar
  \item Archivar, Registrar, Certificar y Operar
  \item Ninguna de las anteriores
\end{enumerate}

\textbf{3. El GDPR es la normativa de protecci\'on de datos de:}
\begin{enumerate}[label=\alph*), leftmargin=2.5em]
  \item Estados Unidos
  \item La Uni\'on Europea
  \item Argentina
  \item Naciones Unidas
\end{enumerate}

\textbf{4. Una cookie de sesi\'on, a diferencia de una persistente:}
\begin{enumerate}[label=\alph*), leftmargin=2.5em]
  \item Se elimina al cerrar el navegador
  \item Dura para siempre
  \item Siempre requiere consentimiento expl\'icito y las persistentes no
  \item No existe, es un mito
\end{enumerate}

\textbf{5. La Ley 26.388 en Argentina:}
\begin{enumerate}[label=\alph*), leftmargin=2.5em]
  \item Regula el teletrabajo
  \item Incorpor\'o los delitos inform\'aticos al C\'odigo Penal
  \item Es la ley de protecci\'on de datos
  \item Regula la firma digital
\end{enumerate}

\textbf{6. ``Privacidad por dise\~no'' significa:}
\begin{enumerate}[label=\alph*), leftmargin=2.5em]
  \item Agregar seguridad al final del desarrollo de un producto
  \item Incorporar la protecci\'on de datos desde el inicio del dise\~no de un producto o servicio
  \item Un estilo de dise\~no gr\'afico minimalista
  \item Una certificaci\'on obligatoria en Argentina
\end{enumerate}

\textbf{7. Antes de aceptar todas las cookies de un sitio, una buena pr\'actica es:}
\begin{enumerate}[label=\alph*), leftmargin=2.5em]
  \item Aceptar siempre para no perder tiempo
  \item Revisar la configuraci\'on y rechazar las no esenciales cuando sea posible
  \item Borrar el navegador
  \item No importa, las cookies no tienen riesgo
\end{enumerate}

\textbf{8. Verdadero o Falso: si una app es gratuita, no recopila ni monetiza datos personales.}
\vspace{4pt}
\checkboxVacio Verdadero \qquad \checkboxVacio Falso

% --- Clave de respuestas ---
\vspace{12pt}
\begin{tcolorbox}[colback=verdeOK!10, colframe=verdeOK,
                  title={\checkMark \; Clave de respuestas}, arc=3mm]
\begin{enumerate}[leftmargin=2em]
  \item \textbf{b)} La protecci\'on de datos personales.
  \item \textbf{a)} Acceder, Rectificar, Cancelar (Suprimir) y Oponerse\ldots
  \item \textbf{b)} La Uni\'on Europea.
  \item \textbf{a)} Se elimina al cerrar el navegador.
  \item \textbf{b)} Incorpor\'o los delitos inform\'aticos al C\'odigo Penal.
  \item \textbf{b)} Incorporar la protecci\'on de datos desde el inicio\ldots
  \item \textbf{b)} Revisar la configuraci\'on y rechazar las no esenciales\ldots
  \item \textbf{Falso.} Muchas apps ``gratuitas'' se financian precisamente a trav\'es de los datos de sus usuarios.
\end{enumerate}
\end{tcolorbox}

\newpage
% ============================================================
%  SECCIÓN 7 — CIERRE DE UNIDAD 3
% ============================================================
\section{Cierre de Unidad 3: el marco normativo y \'etico}

Con esta clase completamos el tablero de la \textbf{Unidad 3}. Ya no solo sabemos c\'omo proteger los sistemas t\'ecnicamente (Unidades 1 y 2), sino que entendemos las reglas legales, \'eticas y sociales bajo las cuales operan los negocios digitales.

\begin{center}
\begin{tikzpicture}[
  node distance=1.2cm and 1.5cm,
  box/.style={rectangle, draw=azulLegal, fill=azulLegal!10,
              text width=3.5cm, align=center, rounded corners=3pt,
              font=\small\bfseries},
  boxsmall/.style={rectangle, draw=moradoEtica, fill=moradoEtica!10,
              text width=3.5cm, align=center, rounded corners=3pt,
              font=\scriptsize},
  arr/.style={-{Stealth[length=2.5mm]}, thick, azulLegal}
]
  \node[box] (ley) {PROTECCI\'ON DE DATOS\\Ley 25.326 / GDPR\\Derechos ARCO};
  \node[box, right=2cm of ley] (delito) {CIBERCRIMEN\\Ley 26.388\\Delitos Inform\'aticos};
  
  \node[boxsmall, below=1.5cm of ley] (priv) {PRIVACIDAD POR DISE\~NO\\\'Etica y Transparencia\\Dark Patterns};
  \node[boxsmall, below=1.5cm of delito] (cook) {COOKIES Y RASTREO\\Consentimiento\\First vs Third Party};

  \draw[arr] (ley.south) -- (priv.north);
  \draw[arr] (delito.south) -- (cook.north);
  \draw[arr, dashed] (priv.east) -- (cook.west);
\end{tikzpicture}
\end{center}

\begin{cajaEtica}[title={El puente hacia la Unidad 4}]
Saber la ley y la \'etica es in\'util si la organizaci\'on no tiene \textbf{reglas internas} para cumplirlas. En la Unidad 4 (Gobierno de la Seguridad) vamos a traducir todo este conocimiento externo en \textbf{Pol\'iticas internas} (como la Pol\'itica de Seguridad Aceptable - PSA) y planes de acci\'on para cuando las cosas salen mal (Gesti\'on de Incidentes).
\end{cajaEtica}

\newpage
% ============================================================
%  SECCIÓN 8 — GLOSARIO Y NOTAS
% ============================================================
\section{Glosario de la Clase 9}

\renewcommand{\arraystretch}{1.4}
\begin{tabularx}{\textwidth}{>{\bfseries}p{4.5cm}X}
\toprule
\textbf{T\'ermino} & \textbf{Definici\'on en una l\'inea} \\
\midrule
Ley 26.388 & Ley argentina que tipifica los delitos inform\'aticos en el C\'odigo Penal. \\
Privacidad por Dise\~no & Integrar la protecci\'on de datos desde la arquitectura inicial del producto. \\
Dark Pattern & Dise\~no de interfaz enga\~noso que manipula las decisiones del usuario. \\
Cookie de Sesi\'on & Archivo temporal que expira al cerrar el navegador. \\
Cookie Persistente & Archivo que permanece en el disco hasta su fecha de expiraci\'on. \\
Cookie de Terceros & Rastreador de un dominio distinto al que el usuario est\'a visitando. \\
Grooming & Delito de contacto digital de un adulto con un menor con fines sexuales. \\
Default Privacy & Privacidad por defecto; el usuario est\'a protegido sin tener que configurar nada. \\
\bottomrule
\end{tabularx}

\vspace{1cm}

\section{Espacio para notas y hallazgos}

\begin{tcolorbox}[colback=grisClaro, colframe=gray!50, arc=2mm, height=7cm]
\textcolor{gray}{\itshape Anot\'a aqu\'i los hallazgos de la Radiograf\'ia de Cookies, ideas para el dise\~no \'etico de tu TP Integrador, o dudas para el Parcial 2.}
\end{tcolorbox}

\vfill
\begin{center}
\small\textcolor{gray}{%
Cuadernillo elaborado para la c\'atedra de Seguridad de la Informaci\'on y de
Sistemas --- Lic.\ en Gesti\'on de Negocios Digitales, UCA.\\
Prohibida su distribuci\'on fuera del \'ambito de la cursada sin autorizaci\'on del docente.}
\end{center}

\end{document}